% !TeX root = ../tfg.tex
% !TeX encoding = utf8

%*******************************************************
% Agradecimientos
%*******************************************************

\section*{Agradecimientos}
\addcontentsline{toc}{section}{Agradecimientos}

\begin{center}
\emph{A mi familia, por haberlo sido todo.}
\end{center}
\vspace{1cm}

\noindent
\textit{Tempus fugit} —el tiempo huye—, decían los antiguos. Y cuánta razón tenían. Aún no soy del todo consciente de lo vertiginosamente rápido que han transcurrido estos cinco años. Comencé la carrera con miedo, con incertidumbre, con ganas de abandonarlo todo y dejarme llevar por la vida, lamentando haber caído en la tentación de matricularme en una titulación tan exigente como un doble grado. \newline

Sin embargo, la paciencia, constancia y valor con los que me han educado mis padres —a quienes dedico especialmente este trabajo— me han permitido, casi sin darme cuenta, llegar al final del TFG y estar a punto de concluir la carrera. A ellos quiero agradecer no solo haber ejercido su papel de padres de forma ejemplar, sino también haberme acompañado fielmente durante todo este proceso y haberme brindado todas las oportunidades para formarme y convertirme en la persona que siempre quise ser. \newline

A vosotros os lo debo todo. Todo lo que soy hoy es fruto de vuestra dedicación en cuerpo y alma a mi crianza y educación. Me siento profundamente afortunado de haber nacido en el seno de esta familia, por el cariño recibido y por la vida tan llena de amor que he tenido. ¿Qué sería del mundo sin el amor? Al fin y al cabo, ser padres no es solo dar la vida —eso sería lo fácil—, sino entregarla, y vosotros lo habéis hecho desde mi concepción.

Eso es lo que más me enorgullece: haber vivido muy feliz con vosotros, mi hermana y mi hermano, sintiéndome profundamente querido día a día, viendo cómo todo lo demás pasa a un segundo plano cuando hay salud, respeto y amor. \newline

Gracias también a toda mi familia, en especial a mis abuelos y a aquellos que ya no están con nosotros. Habéis sido un referente constante de esfuerzo, disciplina, sacrificio, dedicación y esperanza. Aun sin tener demasiados recursos, sin posibilidades de estudiar, viéndoos obligados a trabajar desde edades muy tempranas, lograsteis sacar adelante a vuestra familia —a quienes hoy son mis padres y mis tíos—, y eso es digno de admiración. Sois un ejemplo para mí, y cada día aprendo de vosotros, de vuestra experiencia en la carrera de la vida, de vuestras historias, vivencias y sabiduría. Este trabajo, en parte, también os pertenece. \newline

Quiero agradecer igualmente a la vida por haber puesto en mi camino a mi pareja y a su familia. Gracias por acogerme con tanto cariño, hasta el punto de que, a día de hoy, las líneas entre su familia y la mía están ya muy difuminadas. Hemos recorrido juntos esta ardua etapa, luchando codo con codo contra viento y marea, y llegamos a este momento más unidos y enamorados que nunca, preparados para afrontar una nueva etapa en la vida. \newline

Finalmente, quiero expresar mi más sincero agradecimiento a mi tutor, el profesor José Luis Romero Béjar, por su implicación, cercanía y dedicación con todos los alumnos a los que ha tutorizado, que no son pocos. Gracias por las horas empleadas en revisar este trabajo, por guiarlo con precisión y acierto, por el conocimiento compartido y, en definitiva, por haber contribuido a mi formación como matemático e ingeniero informático. \newline

Ya lo dije allá por tercer curso: hoy en día hay pocas personas que realmente puedan llamarse profesores. Si hay alguien cuya docencia me ha marcado, sin duda ha sido la suya. Basta con observar su entrega y vocación para darse cuenta de que no pasan desapercibidas. \newline

Gracias, de corazón, a todos.

\cleardoublepage

